\documentclass[14pt,a4paper]{extarticle}
\usepackage[T2A]{fontenc}
\usepackage[utf8]{inputenc}
\usepackage[russian]{babel}
\usepackage{geometry}
\usepackage{array,longtable,booktabs}
\usepackage{tikz}
\usetikzlibrary{arrows.meta,positioning}
\usepackage{hyperref}
\geometry{left=30mm,right=15mm,top=20mm,bottom=20mm}
\setlength{\parindent}{1.25cm}
\setlength{\parskip}{0.2em}
\sloppy
\newcommand{\code}[1]{\texttt{\detokenize{#1}}}
\begin{document}
\begin{titlepage}\begin{center}
САНКТ-ПЕТЕРБУРГСКИЙ НАЦИОНАЛЬНЫЙ\\ ИССЛЕДОВАТЕЛЬСКИЙ УНИВЕРСИТЕТ ИТМО
\vspace{1.5cm}

Дисциплина: Бэк-энд разработка
\vspace{1.5cm}

\textbf{Отчет}
\vspace{0.5cm}

Домашняя работа 7
\vspace{0.5cm}

\textbf{Технический дизайн облачной инфраструктуры}
\vfill
\begin{flushright}Выполнил:\\ Цатинян Артём\\ Группа БР1.2
\vspace{0.8cm}

Проверил:\\ Добряков Д. И.\end{flushright}
\vfill Санкт-Петербург\\ 2026 г.
\end{center}\end{titlepage}
\section{Задача}
Разработать технический дизайн облачной инфраструктуры микросервисного приложения бронирования столиков: определить ресурсы, сетевую связность, отказоустойчивость и порядок развёртывания. Документ описывает проектное решение для последующей ЛР7; ресурсы на момент подготовки не создавались.
\section{Исходные данные}
Приложение из ЛР2 реализовано на Kotlin и Spring Boot. Сервис \code{identity-service} отвечает за пользователей и авторизацию, \code{catalog-service} --- за рестораны, меню и сведения о столиках, \code{booking-service} --- за бронирования, \code{review-service} --- за отзывы. Сервисы используют PostgreSQL и принадлежащие им схемы \code{identity}, \code{catalog}, \code{booking}, \code{review}. В ЛР5 всё приложение размещено в Docker Compose на одной виртуальной машине.

Взаимодействие включает синхронные HTTP-вызовы: сервисы проверяют токены через identity, booking обращается к каталогу, а review --- к booking. При создании отзыва review записывает событие в outbox и публикует его через RabbitMQ; catalog получает событие и обновляет агрегат рейтинга. В текущей реализации события обрабатываются с защитой от повторного применения.
\section{Архитектурное решение}
Для учебного стенда предварительно выбран Yandex Cloud: Managed Service for Kubernetes (далее Managed Kubernetes), Managed Service for PostgreSQL, Container Registry и Application Load Balancer. Выбор провайдера и доступные квоты необходимо подтвердить перед ЛР7. Это проектное предположение, а не описание уже созданной инфраструктуры.

Облачный Kubernetes выбран для оркестрации четырёх независимых приложений, управления жизненным циклом контейнеров и последующей настройки HPA. Управляемая PostgreSQL уменьшает объём ручного администрирования. Для учебного варианта брокер RabbitMQ планируется запустить внутри кластера как StatefulSet с постоянным томом; это упрощает состав ресурсов, но оставляет единую точку отказа.

\begin{center}
\begin{tikzpicture}[node distance=7mm, box/.style={draw, rounded corners, align=center, minimum width=10cm, minimum height=8mm}, >=Latex]
\node[box] (client) {Пользователь / клиент API};
\node[box, below=of client] (alb) {Application Load Balancer / Ingress, HTTPS};
\node[box, below=of alb] (ingress) {Ingress-контроллер: маршрутизация по префиксам};
\node[box, below=of ingress] (apps) {Kubernetes: identity, catalog, booking, review};
\node[box, below=of apps] (data) {Managed PostgreSQL \quad | \quad RabbitMQ + постоянный том};
\draw[->] (client)--(alb); \draw[->] (alb)--(ingress); \draw[->] (ingress)--(apps); \draw[->] (apps)--(data);
\end{tikzpicture}

Рисунок 1 --- Логическая схема целевой инфраструктуры
\end{center}
\section{Карта ресурсов}
\begin{longtable}{p{4.0cm}p{5.0cm}p{5.3cm}}
\toprule Компонент & Размещение & Назначение и ограничения \\ \midrule
Четыре приложения & Deployment и ClusterIP Service в Managed Kubernetes & По две реплики в целевом варианте; потребление ресурсов подлежит измерению. \\
HTTP-вход & Application Load Balancer и Ingress & TLS и маршрутизация к сервисам по путям. \\
PostgreSQL & Managed Service for PostgreSQL, один хост для учебного стенда & Одна база с четырьмя схемами; единая точка отказа. \\
RabbitMQ & StatefulSet, Service и постоянный том & Один брокер; при недоступности тома события задерживаются. \\
Образы & Container Registry & Версионирование образов по хешу коммита. \\
Мониторинг & Prometheus, Grafana, Loki, Alloy в кластере & Адаптация конфигурации ЛР6; объём и хранение уточняются при развёртывании. \\
Секреты & Lockbox и Kubernetes Secrets через контролируемую интеграцию & Значения не помещаются в Git или образы. \\
\bottomrule
\end{longtable}
Для PostgreSQL выбран один кластер и отдельные схемы, как в текущем приложении. Это логическая, но не физическая изоляция: отказ хоста затронет все сервисы. Для усиления границ нужны отдельные роли, права на свои схемы и проверка работы Liquibase; переход к отдельным базам требует изменения подключения и миграций.
\section{Сетевая топология}
Планируется VPC с отдельными подсетями для узлов Kubernetes, балансировщика и базы данных в двух доступных зонах. Пример адресного плана: \code{10.20.0.0/24} и \code{10.20.1.0/24} для узлов, \code{10.20.10.0/24} и \code{10.20.11.0/24} для балансировщика, \code{10.20.20.0/24} и \code{10.20.21.0/24} для данных. Конкретные зоны и свободные диапазоны проверяются в аккаунте до создания ресурсов; диапазоны подов и сервисов выбираются без пересечений.

\begin{center}
\begin{tikzpicture}[node distance=7mm, box/.style={draw, rounded corners, align=center, minimum width=11cm, minimum height=8mm}, >=Latex]
\node[box] (pub) {Интернет $\rightarrow$ публичный балансировщик, TCP 443};
\node[box, below=of pub] (worker) {VPC: приватные узлы Kubernetes в двух зонах};
\node[box, below=of worker] (private) {Приватная подсеть данных: Managed PostgreSQL, TCP 5432};
\draw[->] (pub)--(worker); \draw[->] (worker)--(private);
\end{tikzpicture}

Рисунок 2 --- Сетевые уровни приложения
\end{center}
Группы безопасности разрешают входящий HTTPS к балансировщику, служебный и прикладной трафик, необходимый кластеру, и доступ к PostgreSQL только от разрешённых узлов. Правила связи компонентов Kubernetes и health check балансировщика уточняются по документации при внедрении. БД и RabbitMQ не публикуются в интернет. Узлам нужен управляемый исходящий доступ к реестру образов и внешним зависимостям через NAT или доступные приватные подключения; он настраивается отдельно.

Входящие префиксы \code{/identity/}, \code{/catalog/}, \code{/booking/}, \code{/review/} должны направляться на соответствующие сервисы с удалением префикса, как в текущем nginx. Возможность и конкретная настройка переписывания пути выбранным Ingress-контроллером проверяются в ЛР7. Административный доступ к API кластера ограничивается разрешёнными источниками и ролями.
\section{Отказоустойчивость и восстановление}
Для stateless-сервисов планируются две реплики, распределение по зонам, readiness/liveness probes и постепенное обновление. Метрики и пороги проверок подлежат настройке после испытаний. Минимальные значения CPU и памяти будут определены замерами: без них HPA по CPU не имеет корректной базы.

В учебном варианте контрольная плоскость кластера, PostgreSQL и RabbitMQ не заявляются отказоустойчивыми во всех сценариях. Один хост PostgreSQL не обеспечивает высокой доступности; для улучшения целевого варианта требуется конфигурация с несколькими хостами в разных зонах. Managed PostgreSQL поддерживает автоматические резервные копии и восстановление на момент времени; восстановление необходимо проверить на отдельном кластере.

\begin{longtable}{p{3.0cm}p{6.9cm}p{4.4cm}}
\toprule Событие & Ожидаемое поведение & Целевые ориентиры \\ \midrule
Отказ пода приложения & Сервис направляет запросы в готовую реплику, Kubernetes пересоздаёт под. & Восстановление минутного порядка; проверить тестом. \\
Отказ узла & Реплики на другом узле продолжают работу при достаточной ёмкости. & Время зависит от планирования и доступных ресурсов. \\
Отказ зоны с единственной БД или брокером & Данные и обмен событиями недоступны до восстановления соответствующего компонента. & Время простоя заранее не гарантируется. \\
Ошибочное изменение данных & Восстановление кластера БД из резервной копии / на доступный момент времени. & RPO и RTO определяются настройками и проверяются учением. \\
\bottomrule
\end{longtable}
Значения времени восстановления и потери данных пока не измерены и не являются SLA. Доступность приложения при падении БД также не гарантируется: часть запросов зависит от неё синхронно. Для брокера сохранение сообщений зависит от настроек устойчивости очередей, сообщений и тома; это необходимо проверить отдельно.

Для учебного планирования задаются предварительные цели: восстановление прикладного пода за 5 минут, восстановление БД из копии за 120 минут и потеря данных не более 24 часов. Это ориентиры для последующих испытаний, а не подтверждённые характеристики или обещание провайдера. Достижимость целей зависит от настроек резервного копирования, времени последней успешной копии, доступности ресурсов и фактической длительности восстановления; при неудовлетворительном результате архитектуру и цели необходимо пересмотреть.
\section{Безопасность и эксплуатация}
Секреты БД, RabbitMQ и JWT генерируются для окружения и хранятся вне репозитория; учебные пароли из Compose не используются. Внешний трафик защищается TLS; соединение к управляемой БД настраивается с проверкой сертификата. Выдача прав на облачные ресурсы и Kubernetes выполняется по принципу минимально необходимых полномочий. Доступ к Grafana и внутренним интерфейсам следует ограничить.

Контейнеры публикуют логи в стандартные потоки. Для мониторинга предусмотрена адаптация Prometheus и Grafana из ЛР6; для метрик узлов и HPA необходимы соответствующие экспортеры и metrics-server. Проверки OpenAPI из ЛР6 показывают HTTP-доступность, но для readiness в кластере следует добавить отдельные проверки готовности и не полагаться только на документацию API.
\section{План развёртывания}
\begin{enumerate}
\item Подтвердить аккаунт, квоты, доступные зоны и бюджет; оценить стоимость мастера, узлов, хранилищ, БД, балансировщика и сетевого трафика в калькуляторе провайдера.
\item Создать VPC, подсети, маршруты и группы безопасности; затем реестр образов и сервисные учётные записи.
\item Подготовить PostgreSQL, схемы, права и секреты; настроить резервное копирование и TLS.
\item Создать кластер Kubernetes и группы узлов, настроить сетевое подключение, Ingress и доступ к реестру.
\item Подготовить образы четырёх сервисов, перенести конфигурацию подключений, развернуть брокер, затем приложения и балансировщик.
\item Настроить источники метрик и журналов; проверить REST API, сценарий отзыва и очереди, перезапуск пода и восстановление данных.
\item По результатам проверки скорректировать ресурсы, параметры готовности и HPA; документировать порядок очистки ресурсов и резервных копий.
\end{enumerate}
Автоматизация инфраструктуры планируется через Terraform, а приложения --- через Kubernetes-манифесты или Helm. Настоящий документ не утверждает, что эти средства или облачный стенд уже настроены. Плата может начисляться за работающие ресурсы, диски, адреса и другие компоненты; смету следует уточнить до создания стенда.
\section{Вывод}
Подготовлен проект облачной инфраструктуры для четырёх сервисов бронирования ресторанов. Определены размещение компонентов, сетевые границы, варианты восстановления и план реализации. Для ЛР7 остаётся подтвердить провайдера и бюджет и проверить архитектуру на практике.
\section{Источники}
\begin{enumerate}
\item \url{https://yandex.cloud/ru/docs/managed-kubernetes/concepts/high-availability}
\item \url{https://yandex.cloud/ru/docs/managed-kubernetes/operations/connect/security-groups}
\item \url{https://yandex.cloud/ru/docs/managed-postgresql/concepts/backup}
\item \url{https://yandex.cloud/ru/docs/managed-postgresql/concepts/high-availability}
\item \url{https://yandex.cloud/en/docs/managed-kubernetes/concepts/network}
\item \url{https://yandex.cloud/ru/docs/application-load-balancer/concepts/application-load-balancer}
\end{enumerate}
\end{document}
